\chapter{Self-Supervised Learning}
\label{chap:self_supervised}

\begin{tldr}
Self-supervised learning (SSL) learns representations from unlabeled data by defining pretext tasks. This chapter covers contrastive methods (SimCLR, MoCo, CLIP), non-contrastive methods (BYOL, SimSiam, DINO), and masked prediction methods (BERT, MAE). Understanding why these methods work—particularly how they avoid representation collapse—is a key differentiator in staff-level interviews.
\end{tldr}

\section{The Self-Supervised Learning Paradigm}
\label{sec:ssl_paradigm}

\subsection{Motivation}

\textbf{The data labeling bottleneck}:
\begin{itemize}
    \item Labeled data is expensive and slow to collect
    \item Unlabeled data is abundant (text, images, audio everywhere)
    \item Human-provided labels may be noisy or biased
\end{itemize}

\textbf{The SSL solution}: Create supervision signal from the data itself.

\subsection{Taxonomy of SSL Methods}

\begin{figure}[H]
\centering
\begin{tikzpicture}[
    node distance=1.5cm,
    every node/.style={font=\small},
    box/.style={rectangle, draw, rounded corners, minimum width=2.5cm, minimum height=0.8cm, align=center},
    rootbox/.style={box, fill=lightblue, font=\bfseries},
    level1/.style={box, fill=lightgreen},
    level2/.style={box, fill=lightyellow, minimum width=2cm}
]
    % Root
    \node[rootbox] (root) {Self-Supervised Learning};
    
    % Level 1
    \node[level1, below left=1.2cm and 2cm of root] (contrastive) {Contrastive};
    \node[level1, below=1.2cm of root] (noncontr) {Non-Contrastive};
    \node[level1, below right=1.2cm and 2cm of root] (masked) {Masked Prediction};
    
    % Level 2 - Contrastive
    \node[level2, below=0.8cm of contrastive, xshift=-1cm] (simclr) {SimCLR};
    \node[level2, below=0.8cm of contrastive, xshift=1cm] (moco) {MoCo};
    \node[level2, below=1.6cm of contrastive] (clip) {CLIP};
    
    % Level 2 - Non-contrastive
    \node[level2, below=0.8cm of noncontr, xshift=-1cm] (byol) {BYOL};
    \node[level2, below=0.8cm of noncontr, xshift=1cm] (simsiam) {SimSiam};
    \node[level2, below=1.6cm of noncontr] (dino) {DINO};
    
    % Level 2 - Masked
    \node[level2, below=0.8cm of masked, xshift=-0.8cm] (bert) {BERT};
    \node[level2, below=0.8cm of masked, xshift=0.8cm] (mae) {MAE};
    \node[level2, below=1.6cm of masked] (beit) {BEiT};
    
    % Edges
    \draw[-] (root) -- (contrastive);
    \draw[-] (root) -- (noncontr);
    \draw[-] (root) -- (masked);
    
    \draw[-] (contrastive) -- (simclr);
    \draw[-] (contrastive) -- (moco);
    \draw[-] (contrastive) -- (clip);
    
    \draw[-] (noncontr) -- (byol);
    \draw[-] (noncontr) -- (simsiam);
    \draw[-] (noncontr) -- (dino);
    
    \draw[-] (masked) -- (bert);
    \draw[-] (masked) -- (mae);
    \draw[-] (masked) -- (beit);
\end{tikzpicture}
\caption{Taxonomy of self-supervised learning methods.}
\label{fig:ssl_taxonomy}
\end{figure}

\subsection{The Fundamental Challenge: Collapse}

\textbf{Representation Collapse.} When the encoder learns to map all inputs to the same (or very similar) representations, achieving trivially low loss without learning useful features.

\textbf{Types of collapse}:
\begin{enumerate}
    \item \textbf{Complete collapse}: All outputs = constant vector
    \item \textbf{Dimensional collapse}: Outputs span low-dimensional subspace
    \item \textbf{Cluster collapse}: All outputs fall into few clusters
\end{enumerate}

\textbf{Why collapse is tempting}: Many SSL objectives have collapse as a global minimum!

%==============================================================================
\section{Contrastive Learning}
\label{sec:contrastive_learning}
%==============================================================================

\subsection{Core Principle}

\textbf{Key idea}: Learn by contrasting positive pairs (similar) against negative pairs (different).

\begin{figure}[H]
\centering
\begin{tikzpicture}[
    node distance=0.8cm,
    box/.style={rectangle, draw, minimum width=1.5cm, minimum height=0.7cm, rounded corners},
    encoder/.style={box, fill=lightblue},
    aug/.style={box, fill=lightyellow}
]
    % Image
    \node (img) {Image $x$};
    
    % Augmentations
    \node[aug, below left=1cm and 0.5cm of img] (aug1) {Aug $t_1$};
    \node[aug, below right=1cm and 0.5cm of img] (aug2) {Aug $t_2$};
    
    % Augmented images
    \node[below=0.5cm of aug1] (x1) {$x_1$};
    \node[below=0.5cm of aug2] (x2) {$x_2$};
    
    % Encoders
    \node[encoder, below=0.5cm of x1] (enc1) {$f_\theta$};
    \node[encoder, below=0.5cm of x2] (enc2) {$f_\theta$};
    
    % Projectors
    \node[box, fill=lightgreen, below=0.5cm of enc1] (proj1) {$g_\phi$};
    \node[box, fill=lightgreen, below=0.5cm of enc2] (proj2) {$g_\phi$};
    
    % Representations
    \node[below=0.5cm of proj1] (z1) {$z_1$};
    \node[below=0.5cm of proj2] (z2) {$z_2$};
    
    % Contrastive loss
    \node[box, fill=orange!30, below=1cm of $(z1)!0.5!(z2)$] (loss) {Contrastive Loss};
    
    % Positive arrow
    \draw[<->, thick, green!70!black] (z1) -- (z2) node[midway, above] {Positive};
    
    % Edges
    \draw[->] (img) -- (aug1);
    \draw[->] (img) -- (aug2);
    \draw[->] (aug1) -- (x1);
    \draw[->] (aug2) -- (x2);
    \draw[->] (x1) -- (enc1);
    \draw[->] (x2) -- (enc2);
    \draw[->] (enc1) -- (proj1);
    \draw[->] (enc2) -- (proj2);
    \draw[->] (proj1) -- (z1);
    \draw[->] (proj2) -- (z2);
    \draw[->] (z1) -- (loss);
    \draw[->] (z2) -- (loss);
    
    % Negatives (from other images)
    \node[right=2cm of z2, text width=2cm, align=left, font=\small] (neg) {Negatives:\\$z_3, z_4, ..., z_N$\\(other images)};
    \draw[->, thick, red] (neg) -- (loss);
\end{tikzpicture}
\caption{Contrastive learning framework: maximize agreement between augmented views of same image while contrasting with other images.}
\label{fig:contrastive_framework}
\end{figure}

\subsection{SimCLR}

\subsubsection{Algorithm}

\begin{algorithm}[H]
\caption{SimCLR: A Simple Framework for Contrastive Learning}
\begin{algorithmic}[1]
\REQUIRE Batch of images $\{x_k\}_{k=1}^N$, encoder $f$, projector $g$, temperature $\tau$
\FOR{each image $x_k$}
    \STATE Sample two augmentations $t, t' \sim \mathcal{T}$
    \STATE $\tilde{x}_{2k-1} = t(x_k)$, $\tilde{x}_{2k} = t'(x_k)$
\ENDFOR
\STATE $h_i = f(\tilde{x}_i)$ for all $i$ \COMMENT{Representations}
\STATE $z_i = g(h_i)$ for all $i$ \COMMENT{Projections}
\FOR{each positive pair $(i, j)$}
    \STATE $\ell_{i,j} = -\log \frac{\exp(\text{sim}(z_i, z_j)/\tau)}{\sum_{k=1}^{2N} \mathbbm{1}_{k \neq i} \exp(\text{sim}(z_i, z_k)/\tau)}$
\ENDFOR
\STATE $\mathcal{L} = \frac{1}{2N} \sum_{k=1}^{N} [\ell_{2k-1, 2k} + \ell_{2k, 2k-1}]$
\RETURN $\mathcal{L}$
\end{algorithmic}
\end{algorithm}

\subsubsection{Critical Components}

\begin{enumerate}
    \item \textbf{Data Augmentation}
    \begin{itemize}
        \item Random crop and resize (most important)
        \item Color distortion (jitter, grayscale)
        \item Gaussian blur
        \item Composition matters: crop + color best
    \end{itemize}
    
    \item \textbf{Projection Head}
    \begin{itemize}
        \item MLP: Linear → ReLU → Linear
        \item Projects to lower dimension (e.g., 2048 → 128)
        \item Discarded after training; use $h$ not $z$
    \end{itemize}
    
    \item \textbf{Large Batch Size}
    \begin{itemize}
        \item More negatives per positive
        \item Original paper: 4096-8192
        \item Critical for performance
    \end{itemize}
    
    \item \textbf{Temperature}
    \begin{itemize}
        \item Lower $\tau$: Focus on hard negatives
        \item Typical: $\tau = 0.07-0.1$
    \end{itemize}
\end{enumerate}

\begin{keyinsight}[Why Projection Head?]
The projection head allows the encoder $f$ to retain information useful for downstream tasks that might be discarded if optimized directly. The contrastive objective is optimized in a separate space where only invariance matters.

Empirically: Using $h$ (before projection) for downstream tasks outperforms using $z$ (after projection).
\end{keyinsight}

\subsection{MoCo (Momentum Contrast)}

\subsubsection{Motivation}

SimCLR limitation: Needs very large batches for enough negatives.

MoCo solution: Maintain a queue of negatives from past batches.

\subsubsection{Architecture}

\begin{figure}[H]
\centering
\begin{tikzpicture}[
    node distance=0.8cm,
    box/.style={rectangle, draw, minimum width=2cm, minimum height=0.7cm, rounded corners},
    encoder/.style={box, fill=lightblue},
    queue/.style={box, fill=lightyellow, minimum width=3cm}
]
    % Query path
    \node (xq) {$x^q$};
    \node[encoder, below=0.5cm of xq] (encq) {Encoder $f_\theta$};
    \node[below=0.5cm of encq] (q) {$q$};
    
    % Key path
    \node[right=3cm of xq] (xk) {$x^k$};
    \node[encoder, below=0.5cm of xk] (enck) {Encoder $f_{\theta_m}$};
    \node[below=0.5cm of enck] (k) {$k$};
    
    % Queue
    \node[queue, below=1.5cm of $(q)!0.5!(k)$] (queue) {Queue: $\{k_0, k_1, ..., k_K\}$};
    
    % Momentum update
    \draw[->, thick, dashed, orange] (encq.east) -- (enck.west) node[midway, above, font=\small] {momentum};
    
    % Enqueue
    \draw[->, thick] (k) -- (queue) node[midway, right, font=\small] {enqueue};
    
    % Contrastive
    \node[box, fill=lightgreen, below=0.8cm of queue] (loss) {InfoNCE Loss};
    \draw[->] (q) |- (loss);
    \draw[->] (queue) -- (loss);
    
    % Edges
    \draw[->] (xq) -- (encq);
    \draw[->] (xk) -- (enck);
    \draw[->] (encq) -- (q);
    \draw[->] (enck) -- (k);
\end{tikzpicture}
\caption{MoCo architecture with momentum encoder and queue of negatives.}
\label{fig:moco}
\end{figure}

\subsubsection{Key Innovations}

\begin{enumerate}
    \item \textbf{Momentum Encoder}
    \begin{equation}
    \theta_m \leftarrow m \cdot \theta_m + (1 - m) \cdot \theta
    \end{equation}
    Slowly updated ($m = 0.999$) to keep queue representations consistent.
    
    \item \textbf{Queue of Negatives}
    \begin{itemize}
        \item Fixed size (e.g., 65536)
        \item FIFO: Oldest keys dequeued, new keys enqueued
        \item Decouples negative count from batch size
    \end{itemize}
\end{enumerate}

\subsection{CLIP (Contrastive Language-Image Pre-training)}

\subsubsection{Cross-Modal Contrastive Learning}

\begin{figure}[H]
\centering
\begin{tikzpicture}[
    node distance=0.8cm,
    box/.style={rectangle, draw, minimum width=2cm, minimum height=0.7cm, rounded corners},
    encoder/.style={box, fill=lightblue}
]
    % Image side
    \node (img) {Image};
    \node[encoder, below=0.5cm of img] (imgenc) {Image Encoder};
    \node[below=0.5cm of imgenc] (imgz) {$z_I$};
    
    % Text side
    \node[right=4cm of img] (txt) {Text};
    \node[encoder, below=0.5cm of txt] (txtenc) {Text Encoder};
    \node[below=0.5cm of txtenc] (txtz) {$z_T$};
    
    % Similarity matrix
    \node[below=1.5cm of $(imgz)!0.5!(txtz)$] (matrix) {
        \begin{tikzpicture}[scale=0.4]
            \draw[step=0.5] (0,0) grid (3,3);
            % Diagonal
            \foreach \i in {0,...,5} {
                \fill[green!50] (\i*0.5, 2.5-\i*0.5) rectangle (\i*0.5+0.5, 3-\i*0.5);
            }
        \end{tikzpicture}
    };
    \node[below=0.3cm of matrix, font=\small] {Maximize diagonal (matching pairs)};
    
    % Edges
    \draw[->] (img) -- (imgenc);
    \draw[->] (txt) -- (txtenc);
    \draw[->] (imgenc) -- (imgz);
    \draw[->] (txtenc) -- (txtz);
    \draw[->] (imgz) -- (matrix);
    \draw[->] (txtz) -- (matrix);
\end{tikzpicture}
\caption{CLIP learns aligned image-text representations via contrastive learning.}
\label{fig:clip}
\end{figure}

\subsubsection{Loss Function}

\begin{equation}
\mathcal{L} = \frac{1}{2}\left(\mathcal{L}_{I \to T} + \mathcal{L}_{T \to I}\right)
\end{equation}

where:
\begin{align}
\mathcal{L}_{I \to T} &= -\frac{1}{N}\sum_i \log \frac{\exp(z_I^i \cdot z_T^i / \tau)}{\sum_j \exp(z_I^i \cdot z_T^j / \tau)} \\
\mathcal{L}_{T \to I} &= -\frac{1}{N}\sum_i \log \frac{\exp(z_T^i \cdot z_I^i / \tau)}{\sum_j \exp(z_T^i \cdot z_I^j / \tau)}
\end{align}

\subsubsection{Zero-Shot Classification}

CLIP enables zero-shot image classification:
\begin{enumerate}
    \item Create text prompts: ``a photo of a [class]''
    \item Encode prompts to get text embeddings
    \item Encode image to get image embedding
    \item Classify by maximum cosine similarity
\end{enumerate}

\begin{interviewq}
\textbf{Q: What makes a good augmentation policy for contrastive learning, and how does the choice of augmentations define the learned representation?}

\textbf{Difficulty:} L5 \quad \textbf{Type:} Conceptual \quad \textbf{Frequency:} \freqmed{}

\textbf{Strong answer:} Augmentations in contrastive learning define the \textit{invariances} of the learned representation. Two augmented views of the same image are treated as positives, so the encoder learns features that are unchanged across those transformations. A good augmentation policy satisfies three criteria: (1) it removes information irrelevant to the downstream task (e.g., color jitter removes color histogram shortcuts), (2) it preserves semantically meaningful content (e.g., random crop preserves object identity), and (3) it is ``hard enough'' that the model must learn deep features rather than superficial pattern matching.

SimCLR showed that the composition of random crop + color distortion is the single most important combination---either alone is insufficient. Random crop forces the model to match a local region with a global view (learning spatial structure), while color jitter prevents color histogram shortcuts. Gaussian blur adds further difficulty. Critically, augmentation must match the downstream invariance you want: if color is important for your task (e.g., flower species), aggressive color jitter destroys task-relevant information. Multi-crop strategies (DINO uses 2 global + several local crops) improve efficiency by generating more positive pairs per image. The key principle is that augmentations are an implicit form of supervision---they encode your assumptions about what should be invariant.

\textbf{Red flags} (what NOT to say):
\begin{itemize}
    \item ``Just use all augmentations at maximum strength''---this destroys information
    \item Ignoring the relationship between augmentation choice and downstream task
    \item Not knowing that crop + color jitter is the critical combination for SimCLR
\end{itemize}

\textbf{What the interviewer is testing:} Whether the candidate understands that augmentations are not just regularization but define the semantic content of the learned representation.

\textbf{What separates good from great:}
\begin{itemize}
    \item \textbf{L5:} Knows the standard augmentation set and that composition matters
    \item \textbf{L6:} Explains the invariance principle, can reason about task-specific augmentation design
    \item \textbf{L7:} Discusses multi-crop strategies, asymmetric augmentations (DINO), and the information-theoretic view (augmentations control mutual information between views)
\end{itemize}

\textbf{Common follow-ups:} ``What if your downstream task is texture classification---would you change augmentations?'' ``How do multi-crop strategies in DINO differ from SimCLR's two-crop approach?'' ``Can learned augmentations outperform hand-designed ones?''
\end{interviewq}

%==============================================================================
\section{Non-Contrastive Methods}
\label{sec:non_contrastive}
%==============================================================================

\subsection{The Collapse Problem Without Negatives}

Without negatives, what prevents the model from outputting constants?

\textbf{Various solutions}:
\begin{itemize}
    \item BYOL: Momentum encoder + predictor
    \item SimSiam: Stop-gradient
    \item Barlow Twins: Decorrelation loss
    \item VICReg: Variance/invariance/covariance regularization
\end{itemize}

\subsection{BYOL (Bootstrap Your Own Latent)}

\subsubsection{Architecture}

\begin{figure}[H]
\centering
\begin{tikzpicture}[
    node distance=0.6cm,
    box/.style={rectangle, draw, minimum width=1.8cm, minimum height=0.6cm, rounded corners, font=\small},
    encoder/.style={box, fill=lightblue},
    proj/.style={box, fill=lightgreen},
    pred/.style={box, fill=orange!30}
]
    % Online network
    \node (x1) {$x_1$};
    \node[encoder, below=0.4cm of x1] (enc1) {Encoder};
    \node[proj, below=0.4cm of enc1] (proj1) {Projector};
    \node[pred, below=0.4cm of proj1] (pred) {Predictor};
    \node[below=0.4cm of pred] (z1) {$q_\theta(x_1)$};
    
    % Target network
    \node[right=3cm of x1] (x2) {$x_2$};
    \node[encoder, below=0.4cm of x2] (enc2) {Encoder};
    \node[proj, below=0.4cm of enc2] (proj2) {Projector};
    \node[below=0.4cm of proj2] (z2) {$z_\xi(x_2)$};
    
    % Stop gradient
    \node[box, fill=red!20, below=0.4cm of z2] (sg) {stop-grad};
    
    % Loss
    \node[below=1.5cm of $(z1)!0.5!(sg)$] (loss) {$\mathcal{L} = \|q_\theta(x_1) - z_\xi(x_2)\|^2$};
    
    % Labels
    \node[above=0.3cm of x1, font=\bfseries\small] {Online};
    \node[above=0.3cm of x2, font=\bfseries\small] {Target};
    
    % Momentum
    \draw[->, dashed, thick, orange] (enc1.east) -- (enc2.west) node[midway, above, font=\tiny] {momentum};
    \draw[->, dashed, thick, orange] (proj1.east) -- (proj2.west);
    
    % Edges
    \draw[->] (x1) -- (enc1);
    \draw[->] (x2) -- (enc2);
    \draw[->] (enc1) -- (proj1);
    \draw[->] (enc2) -- (proj2);
    \draw[->] (proj1) -- (pred);
    \draw[->] (proj2) -- (z2);
    \draw[->] (pred) -- (z1);
    \draw[->] (z2) -- (sg);
    \draw[->] (z1) -- (loss);
    \draw[->] (sg) -- (loss);
\end{tikzpicture}
\caption{BYOL architecture: Online network predicts target network's representation.}
\label{fig:byol}
\end{figure}

\subsubsection{Why BYOL Doesn't Collapse}

\begin{enumerate}
    \item \textbf{Asymmetry}: Predictor in online network creates architectural asymmetry
    \item \textbf{Momentum target}: Target network provides slowly-moving target
    \item \textbf{Implicit regularization}: Predictor can't easily predict collapsed constant
\end{enumerate}

Recent analysis suggests: The predictor must have limited capacity to prevent learning the identity function.

\subsection{SimSiam}

\subsubsection{Simplification of BYOL}

\begin{itemize}
    \item No momentum encoder (same weights for both branches)
    \item Uses stop-gradient on one branch
    \item Simpler yet effective
\end{itemize}

\begin{equation}
\mathcal{L} = -\frac{1}{2}\left(\frac{p_1}{\|p_1\|} \cdot \text{stopgrad}\left(\frac{z_2}{\|z_2\|}\right) + \frac{p_2}{\|p_2\|} \cdot \text{stopgrad}\left(\frac{z_1}{\|z_1\|}\right)\right)
\end{equation}

\subsubsection{Why Stop-Gradient Prevents Collapse}

\textbf{Why stop-gradient works.} With stop-gradient, the optimization becomes an alternating optimization:
\begin{enumerate}
    \item Fix $z_2$, optimize $p_1$ to match $z_2$
    \item Fix $p_1$, optimize $z_2$ toward the mean of predictions
\end{enumerate}
This converges to non-collapsed equilibria because neither branch can trivially satisfy both objectives simultaneously.

\subsection{DINO (Self-Distillation with No Labels)}

\subsubsection{Key Innovation: Self-Distillation}

\begin{itemize}
    \item Student and teacher see different augmented views
    \item Match output distributions (not just embeddings)
    \item Teacher is momentum-updated student
\end{itemize}

\begin{equation}
\mathcal{L} = -\sum_i p_t^{(i)} \log p_s^{(i)}
\end{equation}

where $p_t, p_s$ are softmax outputs of teacher and student.

\subsubsection{Collapse Prevention}

\begin{enumerate}
    \item \textbf{Centering}: Subtract running mean from teacher outputs
    \begin{equation}
    g_t(x) \leftarrow g_t(x) - c, \quad c \leftarrow mc + (1-m)\bar{g}_t
    \end{equation}
    Prevents single dimension dominating.
    
    \item \textbf{Sharpening}: Apply temperature to teacher
    \begin{equation}
    p_t^{(i)} = \frac{\exp(g_t^{(i)}/\tau_t)}{\sum_j \exp(g_t^{(j)}/\tau_t)}
    \end{equation}
    Lower teacher temperature ($\tau_t < \tau_s$) forces sharp predictions.
\end{enumerate}

\subsubsection{Emergent Properties}

\begin{keyinsight}[DINO Attention Maps]
DINO-trained ViTs learn attention maps that correspond to object segmentation—without any segmentation labels! The [CLS] token attention naturally highlights objects.
\end{keyinsight}

%==============================================================================
\section{Masked Prediction Methods}
\label{sec:masked_prediction}
%==============================================================================

\subsection{Core Idea}

Mask part of the input, predict the masked content from context.

\subsection{BERT: Masked Language Modeling}

(Covered in detail in Chapter~\ref{chap:nlp_architectures})

\textbf{Key points}:
\begin{itemize}
    \item Mask 15\% of tokens
    \item Predict original tokens
    \item 80/10/10 masking strategy
    \item Learns bidirectional representations
\end{itemize}

\subsection{MAE (Masked Autoencoder)}

\subsubsection{Architecture}

\begin{figure}[H]
\centering
\begin{tikzpicture}[
    node distance=1cm,
    decision/.style={diamond, draw, aspect=2.5, inner sep=1pt, fill=yellow!20, font=\small},
    result/.style={rectangle, draw, rounded corners, fill=lightgreen, minimum width=2.5cm, font=\small, align=center},
    arrow/.style={->, thick}
]
    % Start
    \node[decision] (scale) {Corpus size?};
    
    % Scale branches
    \node[decision, below left=1.5cm and 1cm of scale] (small) {Latency?};
    \node[decision, below right=1.5cm and 1cm of scale] (large) {Real-time?};
    
    % Results
    \node[result, below left=1cm and 0cm of small] (cross) {Cross-Encoder};
    \node[result, below right=1cm and 0cm of small] (poly) {Poly-Encoder};
    \node[result, below left=1cm and 0cm of large] (twostage) {Two-Stage\\(Two-Tower + Rerank)};
    \node[result, below right=1cm and 0cm of large] (tower) {Two-Tower Only};
    
    % Edges
    \draw[arrow] (scale) -- node[above left, font=\small] {$<$10K} (small);
    \draw[arrow] (scale) -- node[above right, font=\small] {$>$10K} (large);
    \draw[arrow] (small) -- node[left, font=\small] {Flexible} (cross);
    \draw[arrow] (small) -- node[right, font=\small] {$<$100ms} (poly);
    \draw[arrow] (large) -- node[left, font=\small] {Yes} (twostage);
    \draw[arrow] (large) -- node[right, font=\small] {Batch} (tower);
\end{tikzpicture}
\caption{Decision tree for similarity architecture selection.}
\label{fig:arch_decision}
\end{figure}

\subsubsection{Key Design Choices}

\begin{enumerate}
    \item \textbf{High masking ratio}: 75\% (vs. BERT's 15\%)
    \begin{itemize}
        \item Images have more redundancy than text
        \item Forces learning semantic features, not just interpolation
        \item Makes task sufficiently challenging
    \end{itemize}
    
    \item \textbf{Asymmetric encoder-decoder}
    \begin{itemize}
        \item Encoder: Large (ViT-Large/Huge), processes only visible patches
        \item Decoder: Small (few layers), processes all patches
        \item Efficient: 3-4× speedup over processing all patches
    \end{itemize}
    
    \item \textbf{Pixel reconstruction}
    \begin{itemize}
        \item Predict normalized pixel values
        \item MSE loss on masked patches only
        \item Simpler than predicting discrete tokens (BEiT)
    \end{itemize}
\end{enumerate}

\subsubsection{Why MAE Works}

\begin{keyinsight}
With 75\% masking, the visible 25\% is often insufficient to reconstruct via simple interpolation. The model must learn:
\begin{itemize}
    \item Object structure and shape
    \item Semantic content
    \item Contextual relationships
\end{itemize}
to predict the missing 75\%.
\end{keyinsight}

\subsection{Comparison of Masked Methods}

\begin{table}[H]
\centering
\caption{Comparison of masked prediction methods}
\begin{tabularx}{\textwidth}{lXXXX}
\toprule
\textbf{Method} & \textbf{Domain} & \textbf{Mask Ratio} & \textbf{Target} & \textbf{Key Insight} \\
\midrule
BERT & Text & 15\% & Token IDs & Bidirectional context \\
GPT & Text & N/A (causal) & Next token & Autoregressive \\
MAE & Images & 75\% & Pixels & High masking for vision \\
BEiT & Images & 40\% & Discrete tokens & BERT-style for vision \\
data2vec & Multi-modal & 50\% & Latent features & Unified framework \\
\bottomrule
\end{tabularx}
\end{table}

%==============================================================================
\section{Understanding Why SSL Works}
\label{sec:ssl_theory}
%==============================================================================

\subsection{Information Theory Perspective}

\textbf{InfoNCE Lower Bound.} InfoNCE loss provides a lower bound on mutual information:
\begin{equation}
I(X; Y) \geq \log(K) - \mathcal{L}_{InfoNCE}
\end{equation}
where $K$ is the number of negative samples.

\textbf{Interpretation}: Minimizing InfoNCE maximizes a lower bound on mutual information between different views of the same data.

\subsection{The Augmentation Hypothesis}

\begin{keyinsight}[What Augmentations Define]
Augmentations define what information should be \textbf{invariant} (same across augmentations) vs. \textbf{variant} (changes with augmentation).

Good representations are invariant to augmentation-defined transformations while preserving semantically meaningful information.
\end{keyinsight}

\textbf{Example}: Crop + color jitter
\begin{itemize}
    \item Invariant to: Position, color distribution
    \item Preserves: Object identity, semantic content
\end{itemize}

\subsection{Downstream Transfer}

SSL representations transfer well because:
\begin{enumerate}
    \item They capture semantic structure (needed for pretext task)
    \item They're not overfit to specific labels
    \item Pretext tasks encourage general visual/linguistic understanding
\end{enumerate}

%==============================================================================
\section{Practical Guidelines}
\label{sec:ssl_practical}
%==============================================================================

\subsection{Method Selection}

\begin{table}[H]
\centering
\caption{SSL method selection guide}
\begin{tabularx}{\textwidth}{lXX}
\toprule
\textbf{Scenario} & \textbf{Recommended} & \textbf{Rationale} \\
\midrule
Vision, transfer to classification & MAE, DINO & Strong downstream performance \\
Vision, dense prediction (detection, segmentation) & DINO, DINOv2 & Better local features \\
Vision, limited compute & SimCLR, MoCo & Simpler, well-understood \\
Text understanding & BERT, RoBERTa & Bidirectional context \\
Text generation & GPT-style CLM & Autoregressive natural \\
Vision-language & CLIP & Joint embedding space \\
Multimodal & data2vec, ImageBind & Unified framework \\
\bottomrule
\end{tabularx}
\end{table}

\subsection{Training Tips}

\begin{enumerate}
    \item \textbf{Augmentation is critical}
    \begin{itemize}
        \item Contrastive methods: Strong augmentation essential
        \item MAE: Basic crop/resize sufficient
    \end{itemize}
    
    \item \textbf{Batch size matters for contrastive}
    \begin{itemize}
        \item SimCLR: 4096+ optimal
        \item MoCo: Can work with smaller batches (queue helps)
        \item Non-contrastive: Smaller batches OK
    \end{itemize}
    
    \item \textbf{Training duration}
    \begin{itemize}
        \item SSL typically needs longer training than supervised
        \item 100-1000 epochs common for vision
    \end{itemize}
    
    \item \textbf{Evaluation}
    \begin{itemize}
        \item Linear probe: Freeze encoder, train linear classifier
        \item Fine-tuning: Adapt full model
        \item k-NN: Nearest neighbor in embedding space (no training)
    \end{itemize}
\end{enumerate}

\begin{warningbox}
Augmentation design is the most important decision in contrastive SSL. The augmentations define what information the model considers ``invariant'' vs. ``variant.'' Bad augmentations (e.g., aggressive color jitter for a color-classification task) can destroy task-relevant information and make pre-training counterproductive.
\end{warningbox}

\begin{interviewq}
\textbf{Q: Design an SSL pre-training pipeline for a dataset of 10M unlabeled images and 100K labeled images. Walk through method selection, architecture, training, and evaluation.}

\textbf{Difficulty:} L6 \quad \textbf{Type:} Architecture Design \quad \textbf{Frequency:} \freqhigh{}

\textbf{Strong answer:} I would design a two-stage pipeline: SSL pre-training followed by supervised fine-tuning.

\textit{Stage 1: SSL Pre-training on 10M unlabeled.} Method selection depends on compute budget and downstream task. For classification, I'd choose MAE with ViT-Large---it processes only 25\% of patches (75\% masking), giving 3--4$\times$ training speedup over contrastive methods. For dense prediction (detection, segmentation), I'd use DINOv2 because it produces better local features. If compute is very limited, MoCo v3 with a ResNet-50 is a practical fallback.

\textit{Architecture:} ViT-L/16 encoder, asymmetric decoder (4 layers), pre-train for 800--1600 epochs with cosine LR schedule, AdamW optimizer ($\text{lr}=1.5\text{e-}4$, weight decay 0.05), batch size 4096 across 8--16 GPUs. For MAE, the 75\% masking ratio is well-validated; I'd keep it.

\textit{Stage 2: Fine-tuning on 100K labeled.} Take the pre-trained encoder, add a task-specific head, fine-tune end-to-end with lower learning rate for backbone ($1\text{e-}5$) and higher for head ($1\text{e-}3$). Use 20\% of labeled data as validation. Layer-wise LR decay (0.65 decay factor) helps preserve pre-trained features in early layers.

\textit{Evaluation:} (1) Linear probe accuracy to measure representation quality without fine-tuning confounders, (2) k-NN accuracy as a quick sanity check, (3) full fine-tuning accuracy as the deployment metric. Compare against supervised-only baseline and ImageNet-pretrained baseline.

Expected gain: 5--15\% accuracy improvement over training from scratch with 100K labels, with the gap widening if labeled data is reduced further.

\textbf{Red flags} (what NOT to say):
\begin{itemize}
    \item Jumping to a method without justifying the choice
    \item Forgetting to discuss evaluation strategy for the SSL representations
    \item Not mentioning training duration---SSL requires much longer training than supervised
    \item Ignoring the compute vs.\ quality tradeoff between methods
\end{itemize}

\textbf{What the interviewer is testing:} End-to-end pipeline thinking: method selection, hyperparameter awareness, evaluation strategy, and practical engineering decisions.

\textbf{What separates good from great:}
\begin{itemize}
    \item \textbf{L5:} Knows the pre-train then fine-tune workflow, picks a reasonable method
    \item \textbf{L6:} Justifies method selection based on task and compute, discusses layer-wise LR decay, multi-stage evaluation
    \item \textbf{L7:} Considers domain gap between unlabeled and labeled data, discusses curriculum strategies, addresses when SSL might \textit{not} help (e.g., very different domains), mentions semi-supervised alternatives
\end{itemize}

\textbf{Common follow-ups:} ``What if the unlabeled data is from a different domain?'' ``How would you decide between MAE and DINO?'' ``At what labeled data size does SSL stop being worth the pre-training cost?''
\end{interviewq}

\begin{interviewq}
\textbf{Q: SimCLR vs BYOL vs MAE---give me a decision framework for choosing between them.}

\textbf{Difficulty:} L5 \quad \textbf{Type:} Trade-off \quad \textbf{Frequency:} \freqhigh{}

\textbf{Strong answer:} These three methods represent the three major SSL paradigms---contrastive, non-contrastive, and masked prediction---and the choice depends on four factors: compute budget, batch size constraints, downstream task type, and architecture preference.

\textbf{SimCLR} (contrastive): Needs large batch sizes (4096--8192) because negatives come from the same batch. Requires strong augmentations (crop + color jitter + blur). Best when: you can afford large batches, want a well-understood method, and plan to use CNN backbones. Weakness: performance degrades sharply with small batches; augmentation design is critical. MoCo solves the batch size issue with a momentum queue but adds complexity.

\textbf{BYOL} (non-contrastive): No negatives needed, works with smaller batches (256--512). Uses momentum encoder + predictor for collapse prevention. Best when: you have limited batch size (fewer GPUs), want to avoid augmentation sensitivity, or want simpler training. Weakness: collapse can still occur if hyperparameters are wrong (momentum too low, predictor too large).

\textbf{MAE} (masked prediction): Most compute-efficient (processes only 25\% of patches). Requires ViT architecture. Minimal augmentation needed (just random crop). Best when: you want the best compute-to-quality ratio, plan to use ViT, or need to scale to very large datasets. Weakness: requires many training epochs (1600), less strong for dense prediction than DINO, and the resulting features benefit most from full fine-tuning rather than linear probing.

Quick decision: (1) CNN backbone? $\to$ SimCLR/MoCo or BYOL. (2) ViT backbone, classification? $\to$ MAE. (3) ViT backbone, detection/segmentation? $\to$ DINOv2. (4) Small batch budget? $\to$ BYOL or MoCo.

\textbf{Red flags} (what NOT to say):
\begin{itemize}
    \item ``MAE is always best because it's newest''---it depends on the downstream task
    \item Not knowing that SimCLR requires large batches
    \item Confusing which methods need negative samples
\end{itemize}

\textbf{What the interviewer is testing:} Ability to make principled architecture decisions based on constraints rather than reciting method descriptions.

\textbf{What separates good from great:}
\begin{itemize}
    \item \textbf{L5:} Describes each method correctly and names one tradeoff per method
    \item \textbf{L6:} Provides a structured decision framework, discusses compute and batch size constraints
    \item \textbf{L7:} Mentions linear probe vs.\ fine-tuning gap (MAE is weaker on linear probe), discusses DINOv2 as a unified solution, and considers the data domain
\end{itemize}

\textbf{Common follow-ups:} ``What if you only have 4 GPUs?'' ``Why does MAE have weaker linear probe performance?'' ``How does DINOv2 combine the best of contrastive and masked approaches?''
\end{interviewq}

\begin{interviewq}
\textbf{Q: Your SSL model's downstream performance plateaus despite increasing pre-training data from 1M to 10M images. What could be going wrong?}

\textbf{Difficulty:} L6 \quad \textbf{Type:} Debugging \quad \textbf{Frequency:} \freqmed{}

\textbf{Strong answer:} I would systematically investigate several hypotheses in order of likelihood.

\textit{Hypothesis 1: Augmentation bottleneck.} The augmentations may be too weak, allowing the model to solve the pretext task with shortcuts. Check: are all augmentation components active? Is the task still ``hard'' for the larger dataset? If the model achieves very low contrastive loss early, the augmentations need to be stronger.

\textit{Hypothesis 2: Capacity saturation.} The model may be too small to benefit from more data. A ViT-Base that saturates at 1M images might need scaling to ViT-Large to absorb 10M images' worth of information. Check by plotting linear probe accuracy vs.\ model size at each data scale.

\textit{Hypothesis 3: Training duration.} SSL methods often need hundreds of epochs. With 10$\times$ more data, you may need to maintain the same number of gradient steps (not epochs), meaning 10$\times$ fewer epochs but the same wall-clock compute. If you kept epochs constant, you effectively trained 10$\times$ longer but with diminishing returns.

\textit{Hypothesis 4: Data quality.} The new 9M images may be lower quality, out of domain, or too homogeneous. Duplicate or near-duplicate images inflate dataset size without adding information. Check: run deduplication, verify domain alignment, and evaluate a stratified subset.

\textit{Hypothesis 5: Evaluation mismatch.} The downstream evaluation may have a ceiling. If fine-tuning accuracy is already high, the bottleneck is the downstream task, not the representations. Check linear probe accuracy (ceiling-free) and try harder downstream tasks.

\textbf{Red flags} (what NOT to say):
\begin{itemize}
    \item ``Just train longer''---without diagnosing the root cause
    \item Ignoring data quality as a possible issue
    \item Not considering model capacity as a bottleneck
\end{itemize}

\textbf{What the interviewer is testing:} Systematic debugging of ML pipelines, ability to form and prioritize hypotheses.

\textbf{What separates good from great:}
\begin{itemize}
    \item \textbf{L5:} Mentions data quality and training duration
    \item \textbf{L6:} Provides a structured hypothesis-driven approach, discusses capacity scaling
    \item \textbf{L7:} Connects to scaling laws for SSL (data vs.\ model size vs.\ compute tradeoffs), discusses how to experimentally distinguish hypotheses efficiently
\end{itemize}

\textbf{Common follow-ups:} ``How would you decide between scaling the model vs.\ improving data quality?'' ``What metrics would you track during pre-training to catch this early?'' ``How do you detect near-duplicates at scale?''
\end{interviewq}

\begin{interviewq}
\textbf{Q: BERT vs GPT pre-training objectives---what are the fundamental tradeoffs between masked language modeling and autoregressive language modeling?}

\textbf{Difficulty:} L5 \quad \textbf{Type:} Trade-off \quad \textbf{Frequency:} \freqhigh{}

\textbf{Strong answer:} BERT uses masked language modeling (MLM): mask 15\% of tokens, predict them from bidirectional context. GPT uses causal language modeling (CLM): predict each token from left context only. The tradeoffs are fundamental.

\textit{Context access.} BERT sees both left and right context for every masked position, giving it richer understanding of each token. GPT sees only left context, which is more restrictive but naturally matches the generation task.

\textit{Training signal density.} MLM provides a gradient signal for only 15\% of tokens per forward pass---the masked ones. CLM provides a signal for every token (100\%). This means GPT extracts $\sim$6.7$\times$ more training signal per sequence, making it more data-efficient per token processed.

\textit{Task alignment.} MLM naturally aligns with understanding tasks (classification, NER, QA) because the model learns bidirectional representations. CLM naturally aligns with generation because it directly models $p(x_t | x_{<t})$, which is exactly what we sample during inference. This is why BERT dominates NLU benchmarks while GPT dominates generation.

\textit{Scaling behavior.} CLM scales more favorably because: (1) higher training signal density means better use of compute, (2) the generation interface (prompt $\to$ completion) is universally applicable to any task, and (3) in-context learning emerges at scale. This explains why modern LLMs (GPT-4, LLaMA, Claude) are all decoder-only.

\textit{The train-test mismatch problem.} BERT has a fundamental issue: [MASK] tokens never appear at fine-tuning or inference time, creating distribution shift. The 80/10/10 strategy (80\% mask, 10\% random, 10\% keep) partially mitigates this but doesn't eliminate it. CLM has no such mismatch.

\textbf{Red flags} (what NOT to say):
\begin{itemize}
    \item ``BERT is better for everything because bidirectional is always better''
    \item Not knowing the training signal density difference (15\% vs.\ 100\%)
    \item Ignoring the train-test mismatch from [MASK] tokens
\end{itemize}

\textbf{What the interviewer is testing:} Understanding of how pre-training objectives shape model capabilities and limitations.

\textbf{What separates good from great:}
\begin{itemize}
    \item \textbf{L5:} Knows bidirectional vs.\ causal distinction and typical use cases
    \item \textbf{L6:} Discusses signal density, train-test mismatch, and scaling behavior
    \item \textbf{L7:} Connects to why decoder-only won the scaling race, discusses ELECTRA as an alternative that gives 100\% signal with bidirectional context
\end{itemize}

\textbf{Common follow-ups:} ``Why not just use bidirectional attention in a decoder?'' ``What about T5's span corruption---does it get the best of both?'' ``How does ELECTRA solve the signal density problem?''
\end{interviewq}

\begin{interviewq}
\textbf{Q: Estimate the compute cost of pre-training a BERT-base model on 16GB of text data.}

\textbf{Difficulty:} L6 \quad \textbf{Type:} Estimation \quad \textbf{Frequency:} \freqmed{}

\textbf{Strong answer:} BERT-base has 110M parameters. I'll estimate the total FLOPs and translate to GPU-hours.

\textit{Data size.} 16GB of text at $\sim$4 bytes/token (BPE) $\approx$ 4B tokens. With sequence length 512 and 15\% masking, training processes 4B tokens over $\sim$40 epochs (1M steps $\times$ 256 batch size $\times$ 512 seq length $\approx$ 131B total tokens processed, which is $\sim$33 epochs over 4B tokens).

\textit{FLOPs per token.} The standard approximation is $6 \times P$ FLOPs per token for a forward + backward pass, where $P$ is parameter count. For BERT-base: $6 \times 110\text{M} = 660\text{M}$ FLOPs/token.

\textit{Total FLOPs.} $131\text{B tokens} \times 660\text{M FLOPs/token} \approx 8.6 \times 10^{19}$ FLOPs $\approx 86$ exaFLOPs.

\textit{GPU-hours.} An A100 delivers $\sim$312 TFLOPS (BF16). At $\sim$50\% MFU (Model FLOPs Utilization): effective throughput $\approx$ 156 TFLOPS. Time: $8.6 \times 10^{19} / (156 \times 10^{12}) \approx 5.5 \times 10^5$ seconds $\approx$ 153 GPU-hours. With 8 A100s: $\sim$19 hours wall-clock.

\textit{Sanity check.} The original BERT paper trained on 16 TPUv3 chips for $\sim$4 days. TPUv3 is roughly similar to V100 in FLOPS. 16 $\times$ 96 hours = 1536 GPU-hours, which is higher than my estimate---likely because the original training used FP32 and less efficient implementations. My A100 BF16 estimate is reasonable for modern hardware.

\textbf{Red flags} (what NOT to say):
\begin{itemize}
    \item Not knowing the $6P$ approximation for FLOPs per token
    \item Confusing theoretical peak FLOPS with achievable throughput (forgetting MFU)
    \item Off by more than an order of magnitude on the final estimate
\end{itemize}

\textbf{What the interviewer is testing:} Back-of-envelope estimation skills, understanding of the relationship between model size, data, and compute.

\textbf{What separates good from great:}
\begin{itemize}
    \item \textbf{L5:} Gets the right order of magnitude with reasonable assumptions
    \item \textbf{L6:} Uses the $6P$ rule, accounts for MFU, provides a sanity check
    \item \textbf{L7:} Discusses Chinchilla scaling (BERT is massively over-trained relative to its size), estimates memory requirements, and identifies the batch-size/learning-rate tradeoff
\end{itemize}

\textbf{Common follow-ups:} ``How would the cost change for BERT-large?'' ``Is BERT compute-optimal by Chinchilla standards?'' ``What is the memory footprint during training?''
\end{interviewq}

\begin{interviewq}
\textbf{Q: Next-token prediction vs masked language modeling---why did GPT's autoregressive approach win for generative AI while BERT's bidirectional approach initially dominated NLU benchmarks?}

\textbf{Difficulty:} L6 \quad \textbf{Type:} First Principles \quad \textbf{Frequency:} \freqhigh{}

\textbf{Strong answer:} This is one of the most important questions in modern AI history, and the answer comes down to three interconnected factors: task alignment, scaling properties, and interface universality.

\textit{Task alignment.} Autoregressive models directly model $p(x_t | x_{<t})$, which is exactly what we sample during generation. This means the training objective and the inference procedure are perfectly aligned---no gap. BERT's MLM objective trains the model to fill in blanks, which is a classification task at each masked position. To generate text with BERT, you'd need iterative unmasking, which is unnatural and leads to incoherent outputs. The autoregressive formulation naturally produces coherent, left-to-right text.

\textit{Scaling properties.} Three aspects: (1) Training signal density---CLM uses 100\% of tokens vs.\ MLM's 15\%, making autoregressive training $\sim$6.7$\times$ more efficient per sequence. (2) The loss function provides a natural scaling metric (perplexity) that correlates smoothly with model capability. (3) Emergent abilities---in-context learning, chain-of-thought reasoning---appear in large autoregressive models but not in encoders, because these capabilities require the sequential generation interface.

\textit{Interface universality.} The ``prompt in, completion out'' interface of autoregressive models can express any NLP task: classification (``Is this positive or negative? Answer:''), translation, summarization, reasoning, code generation, and more. BERT's interface is ``fill in [MASK]''---powerful for specific tasks but not a universal interface. This universality means one autoregressive model replaces dozens of task-specific encoder models.

\textit{Why BERT initially won NLU.} Bidirectional context is genuinely better for understanding tasks---seeing both sides of a masked token provides strictly more information. For classification, NER, and QA with modest-sized models (100M--300M params), BERT's architecture is more data-efficient. The shift happened when models became large enough (10B+) that decoder-only models could match encoder performance on NLU tasks while also doing generation.

\textbf{Red flags} (what NOT to say):
\begin{itemize}
    \item ``GPT is just better''---without explaining the structural reasons
    \item Not acknowledging that BERT is still better for specific understanding tasks at smaller scales
    \item Ignoring the training signal density advantage
\end{itemize}

\textbf{What the interviewer is testing:} Deep understanding of why architectural choices matter and how they interact with scale.

\textbf{What separates good from great:}
\begin{itemize}
    \item \textbf{L5:} Explains the generation vs.\ understanding distinction
    \item \textbf{L6:} Discusses signal density, interface universality, and the scale at which decoders catch up on NLU
    \item \textbf{L7:} Connects to the ``bitter lesson'' (scale + simple objective beats clever architecture), discusses why T5 (encoder-decoder with span corruption) was an intermediate step, and mentions recent work on non-causal generative models
\end{itemize}

\textbf{Common follow-ups:} ``Could you design a bidirectional model that generates well?'' ``Why didn't T5's encoder-decoder approach win the scaling race?'' ``What about diffusion-based language models?''
\end{interviewq}

%==============================================================================
\section{Interview Questions}
\label{sec:ssl_interview}
%==============================================================================

\begin{interviewq}
\textbf{Q: Why do non-contrastive methods like BYOL work without negative samples? Shouldn't they collapse to a trivial solution?}

\textbf{Difficulty:} L5 \quad \textbf{Type:} Conceptual \quad \textbf{Frequency:} \freqhigh{}

\textbf{Strong answer:} This is one of the most surprising results in SSL. Without negatives pushing representations apart, the loss landscape has a trivial global minimum: map everything to the same constant vector (zero loss). Yet BYOL avoids this through three complementary mechanisms.

\textit{Architectural asymmetry.} The online network has a predictor head that the target network lacks. This means the online network must \textit{predict} the target representation, not just copy it. If the predictor has limited capacity, it cannot learn the identity function, so collapse becomes a bad local strategy---the predictor would produce a constant output that poorly predicts a moving target.

\textit{Momentum dynamics.} The target network is an exponential moving average (EMA) of the online network ($m = 0.996$). This creates a slowly evolving target. For collapse to occur, both networks would need to converge to the same constant simultaneously, but the momentum dynamics resist this: any perturbation in the online network only slowly propagates to the target.

\textit{Implicit regularization from batch statistics.} BatchNorm (used in the projector/predictor) provides implicit variance regularization across the batch, preventing dimensional collapse where representations span only a low-dimensional subspace.

SimSiam showed that the predictor + stop-gradient alone is sufficient (no momentum needed), framing SSL as alternating optimization: the predictor optimizes to match the current targets, while the encoder optimizes to produce targets that are predictable from augmented views.

\textbf{Red flags} (what NOT to say):
\begin{itemize}
    \item ``It just works because of the momentum encoder''---this misses the predictor's role
    \item Saying negatives are ``implicitly'' present---they are not
    \item Not knowing what happens when you remove the predictor (immediate collapse)
\end{itemize}

\textbf{What the interviewer is testing:} Depth of understanding of collapse prevention---can the candidate go beyond method descriptions to explain the mechanism?

\textbf{What separates good from great:}
\begin{itemize}
    \item \textbf{L5:} Names the three mechanisms (asymmetry, momentum, batch stats)
    \item \textbf{L6:} Explains \textit{why} each mechanism prevents collapse, mentions SimSiam as the simplified case
    \item \textbf{L7:} Discusses the implicit regularization theory, connects to Barlow Twins/VICReg as explicit decorrelation alternatives, and mentions the ongoing theoretical debate
\end{itemize}

\textbf{Common follow-ups:} ``What happens if the predictor is too large?'' ``How does VICReg explicitly prevent collapse?'' ``Is BatchNorm necessary for BYOL to work?''
\end{interviewq}

\begin{interviewq}
\textbf{Q: You have 10M unlabeled images and 10K labeled images in a medical imaging domain. How would you use SSL to improve your classifier, and what challenges do you expect?}

\textbf{Difficulty:} L6 \quad \textbf{Type:} Architecture Design \quad \textbf{Frequency:} \freqhigh{}

\textbf{Strong answer:} Medical imaging adds several domain-specific challenges to the standard SSL pipeline. I would use a multi-stage approach.

\textit{Stage 1: Domain-aware SSL pre-training.} I'd use DINO or DINOv2 with a ViT backbone on the 10M unlabeled medical images. Key adaptations: (1) augmentation design must respect medical image semantics---random crop and flip are fine, but aggressive color jitter may remove diagnostically relevant contrast information; rotation is important since medical images can appear at any orientation. (2) Multi-crop strategy with global-local pairing helps the model learn both holistic and fine-grained features, which is critical for pathology.

\textit{Why DINO over MAE for medical?} DINO's features are better for dense prediction (localization of lesions), and its attention maps naturally segment objects---useful for interpretability in clinical settings. MAE would require more extensive fine-tuning.

\textit{Stage 2: Fine-tuning on 10K labeled.} Given the small labeled set, I'd use a conservative strategy: (1) linear probe first to verify representation quality, (2) then fine-tune with layer-wise LR decay (e.g., 0.65), (3) use aggressive data augmentation during fine-tuning (mixup, cutmix), (4) consider class-balanced sampling since medical datasets often have severe class imbalance.

\textit{Domain challenges:} (1) High inter-image similarity (many normal-looking scans)---augmentations must be strong enough to make the task non-trivial. (2) Class imbalance (rare diseases)---SSL pre-training helps here because the encoder learns general anatomy first. (3) Resolution matters---medical images are often high-resolution; consider using a higher input resolution than standard ImageNet pre-training. (4) Interpretability---clinicians need to understand model decisions; DINO's attention maps provide some interpretability for free.

Expected gain: In medical imaging with limited labels, SSL pre-training typically gives 10--20\% improvement over ImageNet pre-trained models, and 20--40\% over training from scratch.

\textbf{Red flags} (what NOT to say):
\begin{itemize}
    \item Using ImageNet augmentation policies without adaptation for medical domain
    \item Ignoring class imbalance in the labeled dataset
    \item Not considering domain-specific evaluation (sensitivity/specificity, not just accuracy)
\end{itemize}

\textbf{What the interviewer is testing:} Ability to adapt general SSL knowledge to a specific domain with practical constraints.

\textbf{What separates good from great:}
\begin{itemize}
    \item \textbf{L5:} Knows the pre-train then fine-tune workflow, picks a reasonable method
    \item \textbf{L6:} Adapts augmentations for medical domain, discusses class imbalance, considers interpretability
    \item \textbf{L7:} Discusses regulatory constraints (FDA approval requires interpretability), addresses distribution shift across imaging equipment, considers federated SSL for multi-hospital data
\end{itemize}

\textbf{Common follow-ups:} ``Would you use ImageNet pre-trained weights as initialization for SSL?'' ``How do you handle different imaging modalities (CT, MRI, X-ray)?'' ``How would you validate that SSL features capture clinically relevant patterns?''
\end{interviewq}

\begin{interviewq}
\textbf{Q: Explain the InfoNCE loss and its connection to mutual information. Why does temperature matter?}

\textbf{Difficulty:} L6 \quad \textbf{Type:} First Principles \quad \textbf{Frequency:} \freqmed{}

\textbf{Strong answer:} InfoNCE is the core loss function for contrastive SSL. For a positive pair $(i, j)$ with $K$ negatives:

$$\loss_{\text{InfoNCE}} = -\log \frac{\exp(\text{sim}(z_i, z_j)/\tau)}{\sum_{k=1}^{K+1} \exp(\text{sim}(z_i, z_k)/\tau)}$$

\textit{Connection to mutual information.} Oord et al.\ (CPC, 2018) showed that minimizing InfoNCE maximizes a lower bound on mutual information: $I(X; Y) \geq \log(K) - \loss_{\text{InfoNCE}}$. Intuitively, if the model can reliably identify the positive among $K$ negatives, the representations must capture shared information between views. With more negatives, the bound becomes tighter, explaining why larger batch sizes help contrastive methods.

\textit{Why collapse gives non-minimal loss.} If all representations collapse to a constant, every pair has similarity 1, so the loss becomes $\log(K+1)$---the maximum possible value. The minimum is 0 (positive similarity $\gg$ all negative similarities). This is why contrastive methods naturally resist collapse.

\textit{Temperature $\tau$.} Temperature controls the concentration of the softmax distribution. Low $\tau$ ($\sim$0.07): distribution is peaked, gradients are dominated by hard negatives (most similar to the positive). High $\tau$ ($\sim$1.0): distribution is uniform, all negatives contribute equally. The optimal $\tau$ balances two concerns: too low and training is unstable (dominated by a few hard negatives that might be false negatives); too high and the model doesn't learn fine-grained distinctions. In practice, $\tau = 0.07$--$0.1$ works well for vision, while CLIP uses a learnable temperature.

\textbf{Red flags} (what NOT to say):
\begin{itemize}
    \item Confusing InfoNCE with standard cross-entropy---they are related but the denominator structure differs
    \item Not knowing that temperature affects hard negative weighting
    \item Saying ``lower temperature is always better''
\end{itemize}

\textbf{What the interviewer is testing:} Mathematical understanding of the contrastive objective and its information-theoretic motivation.

\textbf{What separates good from great:}
\begin{itemize}
    \item \textbf{L5:} Writes the loss correctly, explains what positive and negative pairs are
    \item \textbf{L6:} Explains the MI lower bound, why collapse gives non-minimal loss, and temperature's role
    \item \textbf{L7:} Discusses tightness of the MI bound, false negatives problem, and alternatives like SigLIP that use sigmoid instead of softmax
\end{itemize}

\textbf{Common follow-ups:} ``What are false negatives in contrastive learning and how do they hurt?'' ``How does SigLIP's loss differ from InfoNCE?'' ``Why does CLIP use a learnable temperature?''
\end{interviewq}

